% Manuscript-ready loss and channel-shift specialization for the journal.
% Uses the draft's existing expectation and theorem notation.
\subsection{A bounded loss satisfying the standing assumptions}
Let $p_w(y\mid x,c)$ denote the softmax confidence assigned to the correct
label under weights $w$ and channel realization $c$. For a fixed
$\varepsilon_{\ell}\in(0,1)$, define the smooth bounded logarithmic loss
\begin{equation}
  \ell_{\varepsilon_{\ell}}(w,c,z)
  = \frac{-\log\!\left[\varepsilon_{\ell}
      +(1-\varepsilon_{\ell})p_w(y\mid x,c)\right]}
    {\log(1/\varepsilon_{\ell})}.
\end{equation}
Since $p_w(y\mid x,c)\in[0,1]$, this loss lies in $[0,1]$ for every
$w,c,z$. Hoeffding's lemma therefore gives, uniformly over the fixed
arguments and for every $\lambda\in\mathbb R$,
\begin{equation}
  \mathbb E\exp\!\left\{\lambda\bigl(L-\mathbb EL\bigr)\right\}
  \leq \exp(\lambda^2/8),
\end{equation}
where $L$ denotes the loss with either the data or the channel variable
randomized as in the standing assumptions. Consequently, the variance
proxies can be chosen as $\sigma^2=\sigma_{\rm tr}^2=1/4$ without
restricting the support of the Gaussian weight posterior or the channel.
For perfect training channels, the loss is deterministic with respect to
the training channel, so the sharper choice $\sigma_{\rm tr}^2=0$ applies.

We use the same bounded loss in the empirical training objective and the
theoretical risk evaluation. Test classification accuracy and ordinary
cross entropy are reported separately; a bound on this normalized loss is
not reported as a numerical bound on unbounded cross entropy or on
classification error. The parameter $\varepsilon_{\ell}=0.01$ is fixed
before the reported campaign. Coefficients $\eta_1,\eta_2$ are tuned
training parameters; reported theoretical penalties are unscaled.

\subsection{A bounded-loss channel-shift specialization}
For any measurable $g:\mathcal C\to[0,1]$ and channel laws $P,Q$,
\begin{equation}
  \left|\mathbb E_Pg(C)-\mathbb E_Qg(C)\right|
  \leq \operatorname{TV}(P,Q).
\end{equation}
This follows by writing
$g(c)=\int_0^1\mathbf1\{g(c)>t\}\,dt$ and applying the defining
event-probability bound for total variation. Applying this inequality to
the bounded loss and averaging over the posterior weights and data gives
a uniform shift penalty independent of the learned posterior. It may
replace the shift term in the risk decomposition; no channel-Lipschitz
constant is needed for this bounded-loss specialization.

For independent coordinate erasures with deployment probability $p$ and
training probability $q$, grouping erasure masks by their number $k$ of
erased coordinates yields
\begin{equation}
  \operatorname{TV}(P_p,P_q)
  =\frac12\sum_{k=0}^{d}\binom dk
  \left|p^k(1-p)^{d-k}-q^k(1-q)^{d-k}\right|.
\end{equation}
In particular, perfect training corresponds to $q=0$, giving
$\operatorname{TV}(P_p,P_0)=1-(1-p)^d$. This remains finite when
$D(P_p\Vert P_0)=\infty$. For perfect training and a nondegenerate
continuous deployment law, total variation equals one, so the resulting
shift guarantee is valid but trivial. When the channel KL is finite,
Pinsker's inequality additionally permits the penalty
$\min\{1,\sqrt{D(P_{\rm ch}\Vert P_{\rm tr})/2}\}$.
